% 05 Conclusion
\section{结论}
\label{sec:conclusion}

本文提出\textbf{渐进式频率条件 FiLM}（Progressive Frequency-conditioned FiLM），一种以 iTransformer 为骨干、通过逐层渐进衰减的 FiLM 将多分辨率频谱信息注入 Transformer Encoder 的轻量高效框架。该方法在 iTransformer 的变量维 token 化范式基础上引入三个核心创新：(1) \textbf{多分辨率频谱提取}——在三个时间窗口上通过 FFT 功率谱提取频带归一化能量、高低频比值与频谱熵，形成尺度鲁棒且信息丰富的频域表征；(2) \textbf{逐层渐进式 FiLM}——压缩后的频谱 latent 通过每层独立的小 MLP 生成调制参数，以指数衰减强度注入每一层 Transformer Encoder，实现由粗到细的层级化频域引导；(3) \textbf{可学习调制强度}——通过 Sigmoid 参数化使模型自适应决定频域信息的注入程度，初始保守设置保证可退化性。

在稳定化设计方面，FiLM 末层零初始化、\(\tanh\) 有界约束与 \(\log(1+x)\) 压缩/截断管道共同确保了训练稳健性与可退化性：初始状态下模型等价于 iTransformer + RevIN，频域调制在训练中逐步激活而非骤然介入。框架通过更换任务头可同时覆盖多变量预测（sequence-to-sequence）与参数回归（sequence-to-vector）两类任务。

实验结果与消融分析（待补充）将进一步验证：(1) 多分辨率频谱相比单窗口频谱的增益；(2) 逐层渐进式调制相比单点调制的优势；(3) 可学习 gamma 相比固定系数的数据集自适应性；(4) 各稳定化组件对训练鲁棒性的贡献。

\subsection{局限性与未来工作}
未来可从以下方向进一步拓展：
\begin{itemize}
  \item 探索可学习的窗口尺寸选择或基于注意力/门控的自适应频谱加权，替代固定三窗口设计；
  \item 引入相位信息或复数频谱特征（利用 FFT 的完整信息），进一步丰富频域表征的完整性；
  \item 在多任务设置下联合预测与异常检测/不确定性估计，以增强工业场景可用性；
  \item 探索频谱 latent \(Z\) 的可视化与解释性，建立频谱熵等指标与预测置信度/难度的量化关联；
  \item 将逐层衰减机制从固定指数衰减推广为可学习的衰减曲线（如轻量 MLP 预测每层衰减系数），进一步降低手工超参依赖；
  \item 研究更严格的理论分析（如频域调制对注意力矩阵谱特性的影响），提升方法在分布外（OOD）条件下的可靠性保障。
\end{itemize}
